\newpage
\begin{center}
    \LARGE
    \textbf{\thesistitle}

    \vspace{0.5cm}

    \authornamefl

    \vspace{1cm}

    \textbf{Abstract}

    \vspace{1cm}
\end{center}

Running services reliably across multiple machines requires solving several difficult problems, such as node discovery, scheduling, failure recovery and access control. Production platforms like Kubernetes solve these problems at the cost of high complexity. This thesis presents the design and implementation of a distributed platform, written from scratch in Go, that provides the essential cluster management mechanisms in a simple and understandable way.

The platform follows a master and worker architecture. Workers are discovered automatically through UDP multicast heartbeats, are grouped into resource pools and run Docker containers with enforced processor and memory limits. The cluster state is replicated across three master nodes using the Raft consensus protocol. Requests sent to a follower are forwarded transparently to the leader, so the failure of a master does not affect users.

The scheduler places the replicas of a service only when the cluster has enough capacity for all of them and chooses the nodes through a pluggable load balancing strategy. Failed nodes are detected within 30 seconds and their containers are restarted on healthy nodes. When a node that was considered dead returns, its duplicate containers are removed automatically. Access to the platform is protected by JWT authentication, role based access control and an audit log, all exposed through a REST API.

The platform was validated through unit tests and manual end to end scenarios covering node discovery, service deployment, fault recovery and security enforcement. The results confirm that the solution fits educational use and smaller deployments.

\vspace{1cm}

\textbf{Keywords:} distributed systems, cluster management, container orchestration, heartbeat, fault tolerance, load balancing, Raft consensus